\subsection{Middleware}\label{subsec:introduction--middleware}

\definition{Middleware} is the \textbf{software layer between the operating system/\break network and the distributed application}.

\begin{figure}[!htp]
    \centering
    \tikzsetnextfilename{middleware}
    \begin{tikzpicture}[
            font=\sffamily,
            app/.style={draw=blue!60!black, fill=blue!10, rounded corners=4pt,
            minimum width=3cm, minimum height=1cm, thick},
            svc/.style={draw=orange!70!black, fill=orange!25, rounded corners=3pt,
            minimum width=2.8cm, minimum height=0.8cm, thick, font=\small},
            plat/.style={draw=gray!70!black, fill=gray!15, rounded corners=4pt,
            minimum width=3cm, minimum height=1cm, thick},
            arr/.style={{Stealth[length=2.5mm]}-{Stealth[length=2.5mm]}, thick, gray!60!black},
            lbl/.style={font=\footnotesize\itshape, text=gray!60!black}
        ]

        % --- Distributed applications ---
        \node[app] (a1) at (-3.6,6.2) {$App_1$};
        \node[app] (a2) at (0,6.2)    {$App_2$};
        \node[app] (a3) at (3.6,6.2)  {$App_3$};
        \node[lbl, rotate=90] at (-5.85,6.2) {distributed applications};

        % --- Middleware band ---
        \draw[thick, orange!70!black, fill=orange!8, rounded corners=6pt]
        (-5.5,2.0) rectangle (5.5,4.6);
        \node[font=\bfseries, text=orange!60!black] at (0,4.2) {Middleware};
        \node[svc] at (-3.6,3.2) {Communication};
        \node[svc] at (0,3.2)    {Naming};
        \node[svc] at (3.6,3.2)  {Coordination};
        \node[lbl] at (0,2.4) {horizontal services shared by all applications};

        % --- Platforms ---
        \node[plat] (p1) at (-3.6,0.4) {Linux};
        \node[plat] (p2) at (0,0.4)    {Windows};
        \node[plat] (p3) at (3.6,0.4)  {Other platform};
        \node[lbl, rotate=90] at (-5.85,0.4) {OS (heterogeneous)};

        % --- Network bar ---
        \draw[thick, gray!70!black, fill=gray!30, rounded corners=4pt]
        (-5.5,-2.0) rectangle (5.5,-1.3);
        \node[font=\small] at (0,-1.65) {Network};

        % --- Arrows ---
        \foreach \x in {-3.6,0,3.6} {
            \draw[arr] (\x,5.7) -- (\x,4.6);
            \draw[arr] (\x,2.0) -- (\x,0.9);
            \draw[arr] (\x,-0.1) -- (\x,-1.3);
        }
        \node[lbl, anchor=west] at (0.15,5.15) {uniform API};
        \node[lbl, anchor=west] at (0.15,1.45) {platform adaptation};

        % --- Takeaway ---
        \node[font=\small\bfseries, text=orange!60!black, text width=10cm, align=center]
        at (0,-2.9) {Middleware hides heterogeneity; it does not eliminate the network.};

    \end{tikzpicture}
    \caption{Middleware as a common layer between distributed applications and heterogeneous platforms.}
\end{figure}

\noindent
The key idea is that applications should not have to deal directly with every low-level difference between machines, operating systems, and network mechanisms.

\highspace
\begin{flushleft}
    \textcolor{Green3}{\faIcon{question-circle} \textbf{Why middleware exists}}
\end{flushleft}
Suppose we have two applications running on different machines: $A$ and $B$. Machine $A$ could use one operating system, while $B$ uses another. They may expose different APIs, use different data representations, or provide different networking facilities.
\begin{itemize}
    \item \textbf{Without middleware}, the application itself has to handle many of these differences:
        \begin{equation*}
            \text{Application}
            \rightarrow
            \text{OS/network-specific details}
        \end{equation*}
        This makes distributed applications harder to build and less portable.

    \item \textbf{Middleware} introduces an intermediate abstraction:
        \begin{equation*}
            \text{Application}
            \rightarrow
            \boxed{\text{Middleware API}}
            \rightarrow
            \text{OS/network}
        \end{equation*}
        So that the application can interact with a \textbf{more uniform interface}.
\end{itemize}
The central purpose is to \hl{hide some platform heterogeneity from distributed applications.}

\highspace
\begin{flushleft}
    \textcolor{Red2}{\faIcon{hourglass-half} \textbf{Network/OS-based distributed system}}
\end{flushleft}
One possibility is to \hl{build the distributed application directly on top of the networking and operating-system facilities}. In this case, the application may need to deal directly with things such as: sockets, network addresses, data encoding, communication details, and differences between platforms. This gives the programmer \hl{a lot of control}, but also exposes \hl{a lot of \textbf{complexity}.}

\highspace
\begin{flushleft}
    \textcolor{Green3}{\faIcon{check-circle} \textbf{Middleware-based distributed system}}
\end{flushleft}
With middleware, we introduce another layer between the application and the operating system/network. The \hl{application uses the services provided by middleware} rather than directly implementing everything using low-level network primitives. This has \textbf{two main purposes}:
\begin{enumerate}
    \item \important{Providing horizontal services}. Middleware provides \textbf{horizontal services} useful to many different distributed applications. ``\emph{Horizontal}''\break means that the \hl{service is not specific to one particular business application.} For example, conceptually, different applications may all need common mechanisms such as communication, naming, and coordination, rather than implementing them independently every time.

        \highspace
        So instead of:
        \begin{equation*}
            App_1 \rightarrow \text{implement everything}
        \end{equation*}

        \begin{equation*}
            App_2 \rightarrow \text{implement everything again}
        \end{equation*}

        we obtain:

        \begin{equation*}
            \begin{array}{ccc}
                App_1 & App_2 & App_3\\
                \downarrow & \downarrow & \downarrow\\
                &\boxed{\text{Middleware}}&
            \end{array}
        \end{equation*}

        The \hl{common distributed functionality is \textbf{reusable}.}

    \item \important{Masking platform differences}. This is the second explicit role of middleware. Imagine a linux machine, a windows machine, and another platform (e.g., MacOS, Android, etc.). At the lower level, these systems may be quite \textbf{different}. \hl{Middleware tries to offer the applications a common abstraction}, while internally adapting to each platform. So middleware acts somewhat like an \textbf{adapter layer}.

        \highspace
        Notice that it does not magically make the underlying machines identical. It simply \textbf{masks some of the differences} so that application developers do not have to handle all of them directly.
\end{enumerate}

\highspace
\begin{examplebox}[: Middleware hides complexity]
    Suppose application $A$ wants to invoke functionality provided by application $B$.
    \begin{itemize}
        \item[\textcolor{Red2}{\faIcon{times}}] Without a middleware abstraction, $A$ might have to think in terms of:
            \begin{equation*}
                \text{IP address}
                +
                \text{socket}
                +
                \text{message format}
                +
                \text{network communication}
            \end{equation*}

        \item[\textcolor{Green3}{\faIcon{check}}] With middleware, the application can instead work with a higher-level operation:
            \begin{equation*}
                A
                \rightarrow
                \boxed{\text{middleware service}}
                \rightarrow
                B
            \end{equation*}
    \end{itemize}
    The low-level communication still happens, of course. Middleware does \textbf{not remove the network}, it just \hl{hides part of its complexity behind a higher-level abstraction.}
\end{examplebox}

\highspace
\begin{takeawaysbox}
    Just three key points to remember about middleware:
    \begin{enumerate}
        \item Middleware sits between applications and the OS/network.
        \item It provides reusable horizontal services for distributed applications.
        \item It masks platform differences, making applications easier to build.
    \end{enumerate}
    So in one sentence: \hl{Middleware is a common software layer that simplifies distributed programming by providing shared services and hiding platform heterogeneity.}
\end{takeawaysbox}
